\chapter{Methods}
\label{ch:methods}

This chapter describes the data, the models, and the two analysis pipelines ---
one quantitative (peer scores), one qualitative (reflective journals) --- that
together answer the four research questions.

\section{Dataset}
\label{sec:dataset}

The data come from COMPSCI~399, the University of Auckland's Computer Science
capstone course. Students work in teams of five or six on a semester-long
software project. Each semester the course runs one or two cohorts, and each
cohort follows the same assessment cycle: four sprints, each ending in a
peer-rating session and a reflective journal submission.

This study draws on six cohorts spanning four academic years
(Table~\ref{tab:cohorts}), covering approximately 174~teams and over 1,000
students. Not every cohort has both peer-session data and journals, and the
LLM-based journal analysis (Section~\ref{sec:llm-pipeline}) was run on four of
the six. Table~\ref{tab:cohorts} records which data sources are available for
each cohort and which analyses each participates in.

\begin{table}[ht]
  \centering
  \caption{Cohort summary. Each cohort runs independently; team labels are
    anonymised and unique within but not across cohorts. ``LLM-coded'' indicates
    the cohort was included in the per-sprint journal analysis
    (Section~\ref{sec:llm-pipeline}).}
  \label{tab:cohorts}
  \begin{tabular}{lrrrcc}
    \toprule
    Cohort        & Students & Teams & Peer sessions & Journals & LLM-coded \\
    \midrule
    2022 Sem~2    &  197     & $\sim$35 & ---        & 5        & ---       \\
    2023 Sem~1    &  110     &  20   &  3            & 4        & ---       \\
    2023 Sem~2    &  219     &  37   &  4            & 5        & \checkmark \\
    2024 Sem~1    &   92     &  15   &  4            & 5        & \checkmark \\
    2024 Sem~2    &  249     &  43   &  4            & 5        & \checkmark \\
    2025 Sem~1    &  142     &  24   &  4            & 5        & \checkmark \\
    \midrule
    Total         & $\sim$1{,}009 & $\sim$174 &      &          &           \\
    \bottomrule
  \end{tabular}
\end{table}

The earliest cohort (2022~Semester~2) has journals but no peer-session data, so
it cannot participate in any analysis that compares journals against peer
scores. The 2023~Semester~1 cohort has both data sources but only three peer
sessions rather than four. The four LLM-coded cohorts (2023~Sem~2 through
2025~Sem~1) form the core of the blind-spot and triage analyses, covering
119~teams and 433 team-sprint observations.

\subsection{Peer-rating data}
\label{sec:peer-data}

At the end of each sprint, every team member distributes a fixed pool of
$10N$~points (where $N$ is the team size) across all members, including
themselves. The course system records the full rating matrix per session. From
this matrix, a student's perceived contribution (PC) is the mean of the peer
scores they received, excluding their self-score. The Individual Weighting
Factor described in Section~\ref{sec:iwf-vulns} is the PC divided by the
expected equal share (10).

Each cohort has three or four peer sessions, yielding a time series of rating
matrices per team. Self-scores are retained in the raw data but excluded from
every analysis in this dissertation, consistent with the course's own
calculation.

\subsection{Reflective journals}
\label{sec:journal-data}

After each sprint, every student submits a reflective journal through the course
management system. The journal has a structured template: students write about
what they did, what the team did, what went well, and what they would change.
Some students write a paragraph; others write a page. The journals are not
graded for content --- they are a completion requirement --- so students have
little incentive to misrepresent, though they may choose what to emphasise.

Each journal maps to a sprint: journal~1 covers the work up to sprint~1,
journal~2 covers sprint~2, and so on. This one-to-one pairing means a team's
journal text for a given sprint can be compared directly with its peer-rating
matrix from the same sprint. The pairing was validated against the peer session
timestamps: the contribution-related flags in the journals (effort imbalance,
under-contribution) agree with low peer scores 85\% of the time.

Across the six cohorts, the journals total over 4,600 individual submissions.
The text is treated as the closest available window into what happened inside a
team --- not as ground truth (students write for assessment), but as an
independent signal that a human reader could use to decide whether a team needs
attention.


\section{Peer Assessment Models}
\label{sec:models}

Four models are applied to each team's rating matrix. All four take the same
input (a matrix of peer scores) and produce the same output (a per-student
weighting factor). They differ in how they handle rater disagreement,
scale~usage, and credibility.

\subsection{Simple Average (baseline)}
\label{sec:simple-avg}

The baseline is the method already used in COMPSCI~399, described in
Section~\ref{sec:iwf-vulns}. Each student's IWF is the mean of the peer scores
they received, divided by the expected equal share. No adjustment is made for
how different raters use the scale, and every rater's scores count equally
regardless of that rater's own standing.

\subsection{WebPA}
\label{sec:webpa-method}

WebPA normalises each rater's scores by dividing by that rater's own mean score
before averaging across raters~\citep{loddington2009webpa}. A student who gives
everyone high marks ends up giving everyone the same normalised score --- their
ratings differentiate nobody. A student who concentrates their points on one
teammate amplifies that differentiation.

Formally, let $s_{ji}$ be the raw score that rater~$j$ gives to student~$i$,
and let $\bar{s}_j$ be rater~$j$'s mean score across all targets. The WebPA
factor for student~$i$ is
\begin{equation}
  W_i \;=\; \frac{1}{|P_i|} \sum_{j \in P_i} \frac{s_{ji}}{\bar{s}_j},
  \label{eq:webpa}
\end{equation}
where $P_i$ is the set of peers who rated student~$i$. The team's mean $W$ is
then used to rescale the factors so they sum to $N$, preserving the team's
overall grade.

\subsection{PeerRank}
\label{sec:peerrank-method}

PeerRank adds a credibility loop: a rater's influence depends on how highly
rated they themselves are~\citep{walsh2014peerrank}. The algorithm iterates.
In round~$t{+}1$, student~$i$'s score is updated as
\begin{equation}
  r_i^{(t+1)} \;=\; (1-\alpha)\,\bar{r}
    \;+\; \alpha \sum_{j \ne i} \frac{r_j^{(t)}}{\sum_{k \ne j} r_k^{(t)}}
    \; s_{ji},
  \label{eq:peerrank}
\end{equation}
where $\bar{r}$ is the team's mean score, $\alpha \in [0,1]$ controls how much
the prior (the mean) is trusted versus the iterative update, and
$r_j^{(t)}/\!\sum_k r_k^{(t)}$ is the normalised credibility of rater~$j$ at
round~$t$. Iteration continues until convergence.

The intuition is simple: if a highly rated student says someone contributed a
lot, that testimony counts more than the same claim from a poorly rated student.
The concern, discussed in Section~\ref{sec:peerrank}, is that this logic was
designed for large networks. In a team of five, one noisy rater is 25\% of all
the evidence, and the credibility loop has very little data to distinguish a
reliable rater from an unreliable one.

\subsection{PeerHITS}
\label{sec:peerhits-method}

PeerHITS separates the two qualities that PeerRank conflates: how much a
student contributed (their \emph{authority} score) and how well they rate others
(their \emph{hub} score). The separation follows HITS~\citep{kleinberg1999hits},
which alternates between two updates:
\begin{align}
  a_i^{(t+1)} &= \sum_{j \ne i} h_j^{(t)} \cdot s_{ji},
  \label{eq:hits-auth} \\
  h_j^{(t+1)} &= \sum_{i \ne j} a_i^{(t)} \cdot s_{ji},
  \label{eq:hits-hub}
\end{align}
where $a_i$ is the authority (contribution) score for student~$i$ and $h_j$ is
the hub (rating quality) score for rater~$j$. After each step, both vectors are
normalised to unit length. A good rater is one whose scores align with the
consensus about who contributed; a good contributor is one whom the good raters
identify.

The practical advantage is that a modest contributor who rates accurately retains
their influence over the weighting. Under PeerRank, that student's low
contribution score would dampen their ratings, losing information. To the
author's knowledge, HITS has not previously been applied to within-team peer
assessment.
